\documentclass[11pt,a4paper]{article}

\usepackage[utf8]{inputenc}
\usepackage[T1]{fontenc}
\usepackage{amsmath,amssymb,amsthm}
\usepackage{mathtools}
\usepackage{booktabs}
\usepackage{tabularx}
\usepackage{hyperref}
\usepackage[margin=2.5cm]{geometry}
\usepackage{natbib}
\usepackage{enumitem}
\usepackage{float}
\usepackage{xcolor}

\newtheorem{theorem}{Theorem}
\newtheorem{proposition}[theorem]{Proposition}
\newtheorem{corollary}[theorem]{Corollary}
\newtheorem{lemma}[theorem]{Lemma}
\newtheorem{definition}[theorem]{Definition}
\newtheorem{remark}[theorem]{Remark}
\newtheorem{example}[theorem]{Example}
\newtheorem{conjecture}[theorem]{Conjecture}
\newtheorem{openquestion}{Open Question}

\newcolumntype{Y}{>{\raggedright\arraybackslash}X}

% -- Colors --
\definecolor{rcolor}{HTML}{D85A30}
\definecolor{icolor}{HTML}{534AB7}
\definecolor{jcolor}{HTML}{1D9E75}
\definecolor{kcolor}{HTML}{185FA5}

\newcommand{\Op}{\mathcal{O}}
\newcommand{\rey}[1]{{\color{rcolor}#1}}
\newcommand{\imag}[1]{{\color{icolor}#1}}
\newcommand{\exac}[1]{{\color{jcolor}#1}}
\newcommand{\recu}[1]{{\color{kcolor}#1}}
\newcommand{\Real}{\operatorname{Re}}
\newcommand{\Imag}{\operatorname{Im}}
\newcommand{\SU}{\mathrm{SU}(2)}

\hypersetup{
  colorlinks=true,
  linkcolor=kcolor,
  citecolor=jcolor,
  urlcolor=rcolor,
  pdftitle={Duality Synthesis in Quaternionic Logic: How Opposites Generate Truth},
  pdfauthor={J. Arturo Ornelas Brand},
  pdfsubject={Duality synthesis, constructive complexity, constructor theory},
  pdfkeywords={duality synthesis, prime algebra, Markov blanket, constructive complexity, superinformation, quaternionic logic, constructor theory}
}

\title{Duality Synthesis in Quaternionic Logic:\\
How Opposites Generate Truth}

\author{J.~Arturo Ornelas Brand\\
  \small Independent Researcher\\
  \small \texttt{arturoornelas62@gmail.com}}
\date{April 13, 2026}

\begin{document}
\maketitle

%======================================================================
\begin{abstract}
%======================================================================
The quaternionic G-lattice defines logical connectives over
$V=[0,1]\times[-1,1]^3$ with an $\SU$ group action that separates
truth ($\Real$) from epistemic character ($\Imag$). The companion
paper establishes that the imaginary dimensions are operationally
distinguishable. This paper shows they are more: through the
synthesis of opposites, they \emph{generate} truth.

We define a synthesis operation on the prime algebra: given dual
concepts $a$ and $\bar{a}$ with prime signatures $\Phi(a)$ and
$\Phi(\bar{a})$, their synthesis
$a\oplus\bar{a}=\mathrm{lcm}(\Phi(a),\Phi(\bar{a}))$ produces a
compound concept whose constructive complexity strictly exceeds that
of either pole. The shared structure $\gcd(\Phi(a),\Phi(\bar{a}))$
forms an algebraic Markov blanket whose interiors are coprime by
the Fundamental Theorem of Arithmetic, yielding a unique partition
that resolves the observer-dependence critique of statistical blankets.

We introduce five axioms (C1--C5) for \emph{constructive complexity
measures} and show that both the prime factor count~$\omega$ and the
assembly index~$MA$ are instances. The gap
operation---extracting what one concept has that another lacks---is
proven to be irreducibly second-order (superinformation). Three
algebraic tasks (directional subsumption, guaranteed composition,
asymmetric gap analysis) are shown to be possible in the prime
medium and impossible under cosine similarity.

The results establish that duality, operating through the imaginary
dimensions of the G-lattice, is a generative mechanism: opposites
do not cancel but produce new truth at higher algebraic layers.
\end{abstract}

%======================================================================
\section{Introduction}
\label{sec:intro}
%======================================================================

The prime-algebraic framework encodes semantic concepts as products
of prime numbers~\citep{p1}, organises them into six algebraic
layers~\citep{p4}, and identifies 14+ candidate dualities across
those layers. The companion paper~\citep{p11} constructs a
quaternionic operator
$\Op=\{0,1,+,\imag{i},\exac{j},\recu{k}\}$ from which all six
layers are recoverable as dimensional projections, and equips the
resulting four-dimensional space with a G-lattice: a residuated
lattice with an $\SU$ group action whose commutation theorem
characterises Boolean logic as the context-free special case.

That paper establishes that the imaginary dimensions
$(\imag{i},\exac{j},\recu{k})$ are operationally distinguishable
from truth ($\rey{r}$): agents at equal truth but different
epistemic profiles make different decisions. Yet the G-lattice's
tautologies depend only on $\Real$; the imaginary dimensions
enrich the \emph{structure around} tautologies but do not alter
which formulas are valid.

This paper addresses the question: \emph{do the imaginary
dimensions produce truth?} We show that they do, through the
synthesis of opposites. When two dual concepts interact, their
\emph{relationship}---encoded in the imaginary dimensions---generates
a compound concept with new crystallised content ($\Real$) that
neither pole possessed individually.

\paragraph{Contributions.}
\begin{enumerate}[nosep]
\item A synthesis operation on the prime algebra
  (Section~\ref{sec:synthesis}) with a proof that it strictly
  increases constructive complexity
  (Theorem~\ref{thm:complexity_increase}).
\item An algebraic Markov blanket theorem: the shared structure
  between any two concepts partitions uniquely into blanket and
  coprime interiors, by FTA
  (Section~\ref{sec:blankets}).
\item Five axioms for constructive complexity measures, with
  $\omega$ and $MA$ as instances
  (Section~\ref{sec:ccm}).
\item A proof that the gap operation constitutes
  superinformation: irreducibly second-order
  (Section~\ref{sec:ct}).
\item Three constructor-theoretic impossibility results
  distinguishing the prime medium from cosine similarity
  (Section~\ref{sec:impossible}).
\item Computational verification on all 14 duality pairs from
  the prime-algebraic ontology, with precise bilateral
  characterisation of which axes increase
  (Section~\ref{sec:computation}).
\item A proof that the prime algebra is irreducible for
  synthesis: no function on quaternionic coordinates alone can
  compute $\mathrm{lcm}$
  (Proposition~\ref{prop:irreducible}).
\item A complexity--constructor bridge theorem: for media with
  unique decomposition, CCM axioms and constructor-theory
  conditions I1--I4 are equivalent; superinformation (I5) is
  the irreducible gap between the frameworks
  (Theorem~\ref{thm:bridge}).
\end{enumerate}

%======================================================================
\section{Preliminaries}
\label{sec:prelim}
%======================================================================

\subsection{The prime algebra}

Each semantic concept $C$ is encoded as a prime signature
$\Phi(C)=\prod_{i}p_i^{s_i(C)}$, where $p_i$ is the $i$-th prime
and $s_i(C)\in\{0,1\}$ indicates whether primitive $i$ is active
for $C$~\citep{p1}. By the Fundamental Theorem of Arithmetic (FTA),
this encoding is injective: distinct concepts have distinct
signatures.

Three operations are defined on signatures:
\begin{itemize}[nosep]
\item $\gcd(\Phi(A),\Phi(B))$: shared structure (what $A$ and $B$
  have in common).
\item $\mathrm{lcm}(\Phi(A),\Phi(B))$: combined structure (the
  smallest concept containing both).
\item $\Phi(A)\bmod\Phi(B)$: subsumption test ($=0$ iff $B$
  divides~$A$).
\end{itemize}

\subsection{The G-lattice}

The quaternionic G-lattice~\citep{p11} is a residuated lattice on
$V=[0,1]\times[-1,1]^3$ equipped with an $\SU$ conjugation action.
The real component $\rey{r}$ encodes crystallised truth; the
imaginary triple $(\imag{i},\exac{j},\recu{k})$ encodes epistemic
character (potentiality, verifiability, recursive selection). The
commutation theorem states that lattice connectives commute with
the group action if and only if they depend exclusively on
$\rey{r}$; Boolean logic is the context-free special case.

\subsection{Dualities}

\begin{definition}[Non-degenerate pair]
\label{def:nondegenerate}
A pair of concepts $(A,B)$ is \emph{non-degenerate} if neither
$\Phi(A)\mid\Phi(B)$ nor $\Phi(B)\mid\Phi(A)$---equivalently,
if both interiors $\iota_A=\Phi(A)/\gcd(\Phi(A),\Phi(B))>1$
and $\iota_B=\Phi(B)/\gcd(\Phi(A),\Phi(B))>1$.
\end{definition}

All theorems in this paper apply to arbitrary non-degenerate
pairs: the synthesis, coprime-interiors, and complexity-increase
results require only that both interiors be non-trivial. The 14
candidate dualities identified in~\citet{p4} are a semantically
motivated \emph{subset} of non-degenerate pairs: they are pairs
whose opposition is grounded in the layer structure (e.g.,
\emph{creaci\'on}/\emph{destrucci\'on},
\emph{orden}/\emph{caos}). Eleven of the 14 are non-degenerate;
three (\emph{vida}/\emph{muerte}, \emph{consciente}/\emph{ausente},
\emph{individual}/\emph{colectivo}) turn out to be degenerate
under dependency expansion---one pole's primitives are a subset
of the other's---so synthesis reduces to the more complex pole
with no strict increase on any axis.

\begin{remark}[Degenerate pairs reflect ontological dependency]
\label{rem:degenerate}
The three degenerate pairs arise from intentional dependency-DAG
design in~\citet{p4}: \emph{muerte} depends on \emph{vida}
(adding destrucci\'on, caos, decaimiento, and muerte itself---19
vs.\ 15~primitives); \emph{ausente} depends on \emph{consciente}
(adding p\'erdida, vac\'io, destrucci\'on---22 vs.\ 18);
\emph{colectivo} depends on \emph{individual} (adding muchos,
cooperaci\'on, uni\'on---20 vs.\ 16). In each case the second pole
\emph{presupposes} the first---death requires life, absence
requires presence, collectivity requires individuality---so
$\Phi(\text{first})\mid\Phi(\text{second})$ by construction. This
is a feature of the ontological design: these dualities express
\emph{existential dependency}, not algebraic opposition. The
synthesis $\mathrm{lcm}$ reduces to the more complex pole because
the ``opposition'' is already resolved within the dependency
structure.
\end{remark}

%======================================================================
\section{The Synthesis Operation}
\label{sec:synthesis}
%======================================================================

\begin{definition}[Synthesis of duals]
\label{def:synthesis}
Given a duality $(a,\bar{a})$ with prime signatures $\Phi(a)$ and
$\Phi(\bar{a})$, the \emph{synthesis} is:
\begin{equation}
a \oplus \bar{a} \;=\; \mathrm{lcm}(\Phi(a),\;\Phi(\bar{a}))
\end{equation}
The synthesis decomposes into three components:
\begin{align}
\text{Blanket:}\quad b &= \gcd(\Phi(a),\;\Phi(\bar{a}))
  \label{eq:blanket}\\
\text{Interior of } a:\quad \iota_a &= \Phi(a)\,/\,b
  \label{eq:int_a}\\
\text{Interior of } \bar{a}:\quad \iota_{\bar{a}} &=
  \Phi(\bar{a})\,/\,b \label{eq:int_abar}\\
\text{Synthesis:}\quad a \oplus \bar{a} &= b \cdot \iota_a \cdot
  \iota_{\bar{a}} \label{eq:synth}
\end{align}
\end{definition}

\begin{example}[Order from opposites]
\label{ex:orden}
In the 72-primitive ontology of~\citet{p4},
\emph{orden} and \emph{caos} each depend on 7~primitives
(including shared ones: \emph{uno}, \emph{fuerza},
\emph{eje\_profundidad}, \emph{m\'as}, \emph{posici\'on\_temporal},
\emph{mover}). The unique primitive of \emph{orden} is
$p_{83}$; that of \emph{caos} is $p_{89}$. Then:
\begin{align*}
b &= \gcd(\Phi(\text{orden}),\,\Phi(\text{caos}))
  &&\text{(6 shared primitives)}\\
\iota_{\text{orden}} &= p_{83}
  &&\text{(the ordering direction)}\\
\iota_{\text{caos}} &= p_{89}
  &&\text{(the disordering direction)}\\
\text{orden}\oplus\text{caos} &= b\cdot p_{83}\cdot p_{89}
  &&(\omega = 6+1+1=8>7)
\end{align*}
The synthesis is not \emph{orden} AND \emph{caos} (that would be
$\gcd = b$, the 6~shared primitives). It is not \emph{orden}
OR \emph{caos} in the lattice sense (component-wise max of
quaternionic coordinates). It is the concept that \emph{contains
both poles and their distinction}---a new entity at a higher level
of constructive complexity. Since both unique primitives are in
Layer~3 ($J$-axis), the bilateral support theorem predicts strict
increase on $\exac{j}$ only, which is confirmed computationally.
\end{example}

\begin{theorem}[Synthesis increases complexity]
\label{thm:complexity_increase}
Let $(a,\bar{a})$ be a duality with $\Phi(a)\neq\Phi(\bar{a})$.
If each pole has at least one prime not shared with the other
(i.e., $\iota_a>1$ and $\iota_{\bar{a}}>1$), then:
\begin{equation}
\omega(a\oplus\bar{a}) > \max(\omega(a),\;\omega(\bar{a}))
\end{equation}
where $\omega(n)$ denotes the number of distinct prime factors of~$n$.
\end{theorem}

\begin{proof}
By the inclusion-exclusion principle for prime factors:
\[
\omega(\mathrm{lcm}(A,B)) = \omega(A) + \omega(B) - \omega(\gcd(A,B))
\]
Let $s = \omega(\gcd)$, $\alpha = \omega(\iota_a) \geq 1$, and
$\beta = \omega(\iota_{\bar{a}}) \geq 1$. Then $\omega(A) = s+\alpha$,
$\omega(B) = s+\beta$, and:
\[
\omega(\mathrm{lcm}) = (s+\alpha)+(s+\beta)-s = s+\alpha+\beta
\]
Since $\alpha\geq1$ and $\beta\geq1$:
\[
s+\alpha+\beta \;\geq\; s+1+\beta \;>\; s+\beta = \omega(B)
\]
and symmetrically $s+\alpha+\beta > s+\alpha = \omega(A)$.
Therefore $\omega(\mathrm{lcm}) > \max(\omega(A),\omega(B))$.
\end{proof}

\begin{remark}[Synthesis is not a lattice connective]
\label{rem:not_connective}
The synthesis $\oplus$ operates on prime signatures, not on
quaternionic coordinates. It is algebraically distinct from the
G-lattice connectives AND ($=\min$) and OR ($=\max$), which
operate component-wise on $V$. The G-lattice connectives preserve
or reduce complexity (min and max do not increase the number of
active components); synthesis \emph{increases} it. This is the
formal sense in which opposites \emph{generate} rather than
\emph{select}.
\end{remark}

\begin{remark}[The synthesis in quaternionic coordinates]
\label{rem:synth_coords}
Theorem~\ref{thm:complexity_increase} operates on prime factor
counts. The Synthesis Embedding Theorem of~\citet{p11} lifts this
to quaternionic coordinates: a layer-based projection
$e\colon\{0,1\}^{72}\to[0,1]^4$ sends each concept's exponent
vector to $(r,i,j,k)$, and satisfies
\[
  e(\mathrm{lcm})_X > \max(e(A)_X,\,e(B)_X)
  \;\;\Longleftrightarrow\;\;
  \mathrm{Int}(A\!\mid\! B)\cap X\neq\emptyset
  \;\text{and}\;
  \mathrm{Int}(B\!\mid\! A)\cap X\neq\emptyset
\]
for each axis $X\in\{R,I,J,K\}$, where
$\mathrm{Int}(A\!\mid\! B)$ is the set of primitives active in $A$
but not in $B$. The strict increase on axis $X$ thus occurs
precisely when both interiors have $X$-support: primitives in the
layer sets that contribute to $X$. This makes the ``synthesis
generates truth'' claim a theorem with an exact characterisation,
not a qualitative observation.
\end{remark}

\begin{proposition}[Irreducibility of the prime algebra]
\label{prop:irreducible}
No function $f\colon[0,1]^4\times[0,1]^4\to[0,1]^4$ computes
synthesis from quaternionic coordinates: there exist concepts
$A,B,C$ with $e(A)=e(B)$ but
$e(\mathrm{lcm}(A,C))\neq e(\mathrm{lcm}(B,C))$.
\end{proposition}

\begin{proof}
The embedding $e$ counts activated primitives per axis without
distinguishing \emph{which} primitives are active. Consider two
concepts $A$ (activating \emph{eje\_vertical}~[Layer~4] and
\emph{mover}~[Layer~3]) and $B$ (activating
\emph{eje\_lateral}~[Layer~4] and
\emph{posici\'on\_temporal}~[Layer~3]). Each has exactly one
Layer-4 and one Layer-3 primitive, so
\[
  e(A) = e(B) = \bigl(\tfrac{1}{36},\;\tfrac{1}{25},\;
  \tfrac{1}{40},\;0\bigr).
\]
Let $C$ activate only \emph{eje\_vertical}. Then
$\Phi(C)\mid\Phi(A)$, so $\mathrm{lcm}(\Phi(A),\Phi(C))=\Phi(A)$
and $e(\mathrm{lcm}(A,C))=e(A)$. But
$\Phi(C)\nmid\Phi(B)$, so $\mathrm{lcm}(\Phi(B),\Phi(C))$ adds
a Layer-4 primitive to~$B$, giving
\[
  e(\mathrm{lcm}(B,C)) = \bigl(\tfrac{2}{36},\;\tfrac{2}{25},\;
  \tfrac{1}{40},\;0\bigr) \;\neq\; e(A).
\]
If $f$ existed, $e(A)=e(B)$ would force
$f(e(A),e(C))=f(e(B),e(C))$, contradicting
$e(\mathrm{lcm}(A,C))\neq e(\mathrm{lcm}(B,C))$.
\end{proof}

\begin{remark}
The irreducibility is a feature of the architecture: the
embedding captures the \emph{structural} aspect of synthesis
(which axes increase, by Remark~\ref{rem:synth_coords}), while
the prime algebra retains the \emph{identity} aspect (which
specific primitives). The two representations are
complementary---the prime algebra is irreducible for computing
synthesis, and the G-lattice is irreducible for characterising
its geometric structure.
\end{remark}

%======================================================================
\section{Algebraic Markov Blankets}
\label{sec:blankets}
%======================================================================

The blanket $b=\gcd(\Phi(a),\Phi(\bar{a}))$ from the synthesis
decomposition has a structural property that resolves a standing
critique of statistical Markov blankets.

\begin{theorem}[Coprime interiors]
\label{thm:mb}
For any two concepts $C_1,C_2$ with prime signatures
$\Phi(C_1),\Phi(C_2)$, define:
\begin{itemize}[nosep]
\item Blanket: $b = \gcd(\Phi(C_1),\Phi(C_2))$
\item Interior of $C_1$: $\iota_1 = \Phi(C_1)/b$
\item Interior of $C_2$: $\iota_2 = \Phi(C_2)/b$
\end{itemize}
Then $\gcd(\iota_1,\iota_2) = 1$.
\end{theorem}

\begin{proof}
By FTA, write $\Phi(C_j) = \prod_i p_i^{a_i^{(j)}}$. Then:
\begin{equation}
b = \prod_i p_i^{\min(a_i^{(1)}, a_i^{(2)})}, \qquad
\iota_j = \prod_i p_i^{a_i^{(j)} - \min(a_i^{(1)}, a_i^{(2)})}
\end{equation}
For each prime $p_i$, let
$e_i = a_i^{(1)} - \min(a_i^{(1)}, a_i^{(2)})$ and
$f_i = a_i^{(2)} - \min(a_i^{(1)}, a_i^{(2)})$.

If $a_i^{(1)} \leq a_i^{(2)}$, then $e_i=0$, so $p_i\nmid\iota_1$.
If $a_i^{(1)} > a_i^{(2)}$, then $f_i=0$, so $p_i\nmid\iota_2$.
In both cases $\min(e_i,f_i)=0$. Therefore:
\[
\gcd(\iota_1,\iota_2) = \prod_i p_i^{\min(e_i,f_i)} = 1 \qedhere
\]
\end{proof}

\begin{remark}[Algebraic Markov property]
The condition $\gcd(\iota_1,\iota_2)=1$ is the algebraic analogue
of conditional independence $\iota_1\perp\iota_2\mid b$: once the
shared structure (blanket) is factored out, the remainders
(interiors) share \emph{nothing}. No prime factor ``leaks'' from one
interior to the other through the blanket.
\end{remark}

\begin{corollary}[Unique partition]
\label{cor:unique}
The partition $(\iota_1,b,\iota_2)$ is unique for any pair
$(C_1,C_2)$.
\end{corollary}

\begin{proof}
By FTA, $\Phi(C_1)$ and $\Phi(C_2)$ have unique prime
factorisations. Hence $b=\gcd(\Phi(C_1),\Phi(C_2))$ and
$\iota_j=\Phi(C_j)/b$ are each uniquely determined.
\end{proof}

\begin{remark}[Relation to statistical Markov blankets]
The prime-algebraic blanket addresses the \emph{uniqueness} component
of the observer-dependence critique~\citep{bruineberg2018}: FTA
guarantees a single partition, discovered by factorisation rather
than chosen by an observer. The conditional-independence analogy
($\gcd(\iota_1,\iota_2)=1$ as ``no shared prime leaks through
the blanket'') is structural, not probabilistic: the prime algebra
has no probability measure, so conditional independence in the
statistical sense does not apply. What is gained is an exact,
observer-free partition; what is not claimed is a reduction to
the probabilistic framework.
\end{remark}

\begin{proposition}[Nested blankets]
\label{prop:nested}
If $\Phi(A)\mid\Phi(B)\mid\Phi(C)$ (divisibility chain), then
at each level the Markov property holds: the interior of $A$
relative to $B$ and the interior of $B$ relative to $C$ are
coprime.
\end{proposition}

\begin{proof}
Apply Theorem~\ref{thm:mb} at each level of the chain. The FTA
argument is identical at each scale.
\end{proof}

%======================================================================
\section{Constructive Complexity Measures}
\label{sec:ccm}
%======================================================================

The synthesis operation increases the prime factor count $\omega$.
We now show that $\omega$ belongs to a general class of measures
that detect constructive depth.

\begin{definition}[Constructive complexity measure]
\label{def:ccm}
A \textbf{constructive complexity measure} over a domain $D$ with
building blocks $B$ and composition $\oplus$ is a function
$\mu:D\to\mathbb{Z}_{\geq0}$ satisfying:
\begin{enumerate}[nosep]
\item[(C1)] $\mu(b)=0$ for all $b\in B$ \hfill(grounding)
\item[(C2)] $\mu(x)\geq0$ for all $x\in D$ \hfill(non-negativity)
\item[(C3)] $\mu(a\oplus b)\leq\mu(a)+\mu(b)+c$ for constant
  $c\geq0$ \hfill(subadditivity)
\item[(C4)] $\mu$ is computable \hfill(computability)
\item[(C5)] For all $a\in D$ and $b\in B$ with $b\nmid a$:
  $\mu(a\oplus b)>\mu(a)$ \hfill(strict monotonicity)
\end{enumerate}
\end{definition}

\begin{remark}
C5 excludes degenerate measures (e.g., $\mu\equiv0$) and ensures
that adding genuinely new building blocks always increases
complexity.
\end{remark}

\subsection{Instance 1: Assembly index}

Walker \& Cronin's assembly index $MA$~\citep{sharma2023} operates
on molecular structures with chemical joining as composition and
$c=1$. It counts the minimum number of joining operations needed
to construct an object from basic fragments. High $MA$ combined
with high copy number is proposed as a substrate-independent
signature of selection.

\subsection{Instance 2: Prime factor count}

The join count $\mu(x)=\omega(x)-1$ (for $x>1$; $\mu(1)=0$)
operates on semantic concepts with $\mathrm{lcm}$ as composition
and $c=1$. Here $\omega$ is the number of distinct prime factors.
C1--C5 are immediate: grounding ($\omega(p_i)=1$, so $\mu=0$),
non-negativity, subadditivity (by inclusion-exclusion),
computability (by trial division), and strict monotonicity (adding
a new prime increases $\omega$ by~1).

\subsection{Structural parallel}

\begin{proposition}[Structural parallel]
\label{prop:parallel}
$MA$ and $\omega$ share three properties beyond C1--C5:
\begin{enumerate}[nosep]
\item \textbf{Overlap correction:}
  $\omega(\mathrm{lcm}(A,B))=\omega(A)+\omega(B)-\omega(\gcd(A,B))$;
  for $MA$, shared sub-assemblies reduce the combined index
  analogously.
\item \textbf{Finite irreducible basis:} primes (resp.\ basic
  molecular fragments).
\item \textbf{Discrimination of structured from random:} the Domain
  Validity Score~\citep{p4} perfectly discriminates 8 real domains
  from 3 pseudoscience controls (gap $0.233$, zero errors).
\end{enumerate}
They diverge on one key point: $\omega$ has \textbf{unique
decomposition} (FTA), while $MA$ does not---multiple assembly
pathways may yield the same index. This makes $\omega$ strictly
stronger as an algebraic invariant.
\end{proposition}

\begin{proposition}[Rarity of high complexity under random assembly]
\label{prop:rarity}
In a prime product space with $k$ building blocks, under uniform
random selection from all $2^k$ possible objects, the probability
of obtaining an object with complexity $\mu\geq m$ is the upper
tail of $\mathrm{Binomial}(k,1/2)$:
\begin{equation}
P(\mu\geq m) = \frac{1}{2^k}\sum_{j=m+1}^{k}\binom{k}{j}
\end{equation}
For $m\gg k/2$, this decays exponentially:
$P(\mu\geq m)\leq\exp(-2(m-k/2)^2/k)$ by Hoeffding's inequality.
\end{proposition}

\begin{proof}
The number of objects with exactly $j$ prime factors is
$\binom{k}{j}$. Under uniform selection,
$P(\omega=j)=\binom{k}{j}/2^k = \mathrm{Binomial}(k,1/2)$.
Since $\mu=\omega-1$, $P(\mu\geq m)=P(\omega\geq m+1)$.
The exponential bound follows from Hoeffding's inequality.
\end{proof}

\begin{conjecture}[Universal Selection Signature]
\label{conj:uss}
For any constructive complexity measure $\mu$ satisfying C1--C5
over a domain $D$ with $k$ building blocks: the empirical
distribution of $\mu$ values in a population produced by
selection is stochastically dominated by the Binomial$(k,1/2)$
null of Proposition~\ref{prop:rarity}, with separation detectable
by a one-sided Kolmogorov--Smirnov test at sample sizes
polynomial in~$k$.
\end{conjecture}

\begin{remark}[Testability]
The USS is not straightforwardly testable on the prime-algebraic
ontology: the 72~primitives \emph{are} the building blocks, so
their $\omega$ values ($1$--$22$, mean~$9.4$) are low by
construction, not by selection. The relevant null (uniform over
$2^{72}$ products, mean~$\omega=36$) applies to \emph{compound
concepts}, not to the basis. The closest existing empirical test
is the Domain Validity Score discrimination of~\citet{p4} (gap
$0.233$, zero misclassifications across 11~domains), which
confirms that $\omega$-derived measures detect structure in real
data. A direct USS test requires a corpus of naturally-occurring
compound concepts with known prime signatures.
\end{remark}

%======================================================================
\section{The Prime Algebra as Information Medium}
\label{sec:ct}
%======================================================================

We verify that the prime algebra satisfies the five conditions for
an information medium under Deutsch \& Marletto's Constructor
Theory of Information~\citep{deutsch2015}.

\begin{proposition}[Conditions I1--I4]
\label{prop:ct}
The prime algebra satisfies:
\begin{enumerate}[nosep]
\item[\textbf{I1}] \textbf{Distinguishability:} For distinct
  $C_1\neq C_2$, $\Phi(C_1)\neq\Phi(C_2)$ by construction.
  The factorisation algorithm recovers the unique decomposition.
\item[\textbf{I2}] \textbf{Clonability:} Integer assignment copies
  any signature exactly to a blank substrate (memory initialised
  to $1$).
\item[\textbf{I3}] \textbf{Non-triviality:} With $k$ primes,
  $2^k\geq2$ distinguishable states.
\item[\textbf{I4}] \textbf{Interoperability:} Measurement
  (factorisation), copying (assignment), and computation
  ($\bmod$, $\mathrm{lcm}$, $\gcd$) are all exact, deterministic,
  and composable within the same integer medium.
\end{enumerate}
\end{proposition}

\subsection{Superinformation (Condition I5)}

\begin{proposition}[Condition I5 satisfied]
\label{prop:super}
The gap operation
$\mathrm{gap}(A,B)=\Phi(A)/\gcd(\Phi(A),\Phi(B))$
constitutes superinformation in the prime-algebraic medium.
\end{proposition}

\begin{proof}
Three requirements.

\textbf{(i) The gap is second-order.} A first-order information
task maps a single variable to an output:
$f:\Phi(A)\mapsto\Phi(A')$. The gap requires \emph{two} variables
as input and outputs a \emph{relationship}---the structural
difference between two concepts---not the content of either alone.

\textbf{(ii) The gap is irreducible to first-order.} Suppose a
function $g:\mathbb{Z}_{>0}\to\mathbb{Z}_{>0}$ exists with
$\mathrm{gap}(A,B)=g(\Phi(A))$ for all $B$. Fix $A$ with
$\Phi(A)=p_1 p_2 p_3$ and choose $B_1$ with $\Phi(B_1)=p_1$,
$B_2$ with $\Phi(B_2)=p_1 p_2$. Then
$\mathrm{gap}(A,B_1)=p_2 p_3 \neq p_3=\mathrm{gap}(A,B_2)$,
but $g(\Phi(A))$ is a fixed value. Contradiction.

\textbf{(iii) The output remains in the medium.} Since
$\gcd(\Phi(A),\Phi(B))$ divides $\Phi(A)$, the quotient is a
positive integer whose prime factors are a subset of those
of~$\Phi(A)$.
\end{proof}

\begin{remark}[Asymmetry]
In general $\mathrm{gap}(A,B)\neq\mathrm{gap}(B,A)$: the gap
encodes an \emph{oriented} relationship. This directional structure
is characteristic of superinformation.
\end{remark}

\begin{remark}[Gap as the engine of synthesis]
The gap is exactly the interior from
Definition~\ref{def:synthesis}: $\mathrm{gap}(a,\bar{a})=\iota_a$
and $\mathrm{gap}(\bar{a},a)=\iota_{\bar{a}}$. The
superinformation proof therefore shows that the
\emph{distinguishing content} of each pole---the raw material of
synthesis---is irreducibly second-order. Synthesis is not reducible
to operations on either pole alone; it requires the relationship
between them.
\end{remark}

%======================================================================
\section{The Possible/Impossible Distinction}
\label{sec:impossible}
%======================================================================

The prime algebra provides a concrete instance of Constructor
Theory's core innovation: a distinction between possible and
impossible tasks across information media.

\begin{theorem}[Three impossibility results]
\label{thm:impossible}
The following three tasks are \textbf{possible} in the
prime-algebraic medium and \textbf{impossible} under cosine
similarity:
\end{theorem}

\begin{proof}
\textbf{Task 1: Directional subsumption.}
$A\supseteq B$ is directional: $A\supseteq B$ does not imply
$B\supseteq A$.
\begin{itemize}[nosep]
\item \textbf{Prime:} $\Phi(A)\bmod\Phi(B)=0\iff A\supseteq B$.
  Directional by construction. \textit{Possible.}
\item \textbf{Cosine:} $\cos(\vec{A},\vec{B})=
  \cos(\vec{B},\vec{A})$ (inner product is symmetric). No function
  of cosine alone can distinguish ``$A$ contains $B$'' from
  ``$B$ contains $A$.'' \textit{Impossible.}
\end{itemize}

\textbf{Task 2: Guaranteed composition.}
Produce $C$ guaranteed to subsume both $A$ and $B$.
\begin{itemize}[nosep]
\item \textbf{Prime:} $C=\mathrm{lcm}(\Phi(A),\Phi(B))$.
  Algebraically guaranteed. \textit{Possible.}
\item \textbf{Cosine:} $\vec{C}=\vec{A}+\vec{B}$ does not
  guarantee $\cos(\vec{C},\vec{A})\geq\tau$ for any threshold
  (e.g., $\vec{A}=-\vec{B}$ gives $\vec{C}=\vec{0}$).
  \textit{Impossible.}
\end{itemize}

\textbf{Task 3: Asymmetric gap analysis.}
Extract what $A$ has that $B$ doesn't, and vice versa, as
distinct objects.
\begin{itemize}[nosep]
\item \textbf{Prime:}
  $\iota_A=\Phi(A)/\gcd(\Phi(A),\Phi(B))$,
  $\iota_B=\Phi(B)/\gcd(\Phi(A),\Phi(B))$, with
  $\gcd(\iota_A,\iota_B)=1$ by FTA. \textit{Possible.}
\item \textbf{Cosine:} $\vec{A}-\vec{B}=-(\vec{B}-\vec{A})$ is
  antisymmetric but does not decompose into separable components
  unique to each side. \textit{Impossible.}
\end{itemize}
\end{proof}

\begin{remark}
Task~1 (directional subsumption) is impossible under \emph{any}
symmetric similarity, not only cosine: the inner product
$\langle\vec{A},\vec{B}\rangle=\langle\vec{B},\vec{A}\rangle$
cannot distinguish ``$A$ contains $B$'' from ``$B$ contains $A$.''
Tasks~2 and~3 are impossible under cosine specifically---the
standard metric in modern NLP (word2vec, BERT, sentence
transformers)---but could in principle be recovered by richer
vector-space operations (e.g., element-wise max for Task~2,
signed residuals for Task~3). The practical force of the
impossibility is that cosine is what practitioners use.
\end{remark}

\begin{remark}
The three tasks map directly to the synthesis decomposition:
Task~1 tests subsumption ($\Phi(a)\mid\Phi(c)$), Task~2 computes
synthesis ($\mathrm{lcm}$), and Task~3 extracts the interiors
($\iota_a,\iota_{\bar{a}}$). The impossibility results therefore
show that the \emph{synthesis of opposites} is impossible in flat
vector spaces---it requires the Boolean lattice structure of the
prime algebra.
\end{remark}

%======================================================================
\section{The Rule of Three as Constructor}
\label{sec:rule3}
%======================================================================

The synthesis $a\oplus\bar{a}=\mathrm{lcm}(\Phi(a),\Phi(\bar{a}))$
combines two concepts. A related operation---analogical
completion---requires three. The Rule of Three,
$C_4=(\Phi(C_2)\times\Phi(C_3))/\Phi(C_1)$, introduced
in~\citet{p1}, solves $A:B=C:?$ deterministically when
$\Phi(C_1)\mid\Phi(C_2)\times\Phi(C_3)$.

This satisfies Deutsch's definition of a
\textbf{constructor}~\citep{deutsch2013}: an entity that causes a
specific transformation and retains the capacity to repeat it. It
also constitutes a fourth impossibility under cosine similarity:
the proportional relationship $\Phi(C_4)/\Phi(C_3)=\Phi(C_2)/\Phi(C_1)$
requires directed ratios, not symmetric dot products.

\begin{remark}[Relation to synthesis]
The Rule of Three and the synthesis operation are complementary
constructors on the prime algebra. Synthesis combines
\emph{opposites} ($\mathrm{lcm}$: what both have plus what
distinguishes them). The Rule of Three completes
\emph{proportions} (multiplication and division: the structural
transfer from one pair to another). Both are possible in the
prime medium and impossible under cosine similarity; both are
deterministic constructors in Deutsch's sense.
\end{remark}

%======================================================================
\section{Computational Verification}
\label{sec:computation}
%======================================================================

All results were verified on the 72 primitives and 14 duality
pairs of~\citet{p4}. The layer-based embedding of~\citet{p11}
maps each concept's dependency-expanded signature to
$(r,i,j,k)\in[0,1]^4$, with axis-contributing sets
$|R|=36$, $|I|=25$, $|J|=40$, $|K|=4$.

\begin{table}[H]
\centering
\caption{Verification of synthesis properties on the 14 duality
pairs from~\citet{p4}.}
\label{tab:verification}
\begin{tabularx}{\textwidth}{@{}lrrl@{}}
\toprule
\textbf{Test} & \textbf{N} & \textbf{Violations} & \textbf{Result} \\
\midrule
$\omega(\mathrm{lcm})>\max(\omega(a),\omega(\bar{a}))$
  for non-degenerate pairs & 11 & 0 & exact \\
Coprime interiors ($\gcd(\iota_a,\iota_{\bar{a}})=1$)
  & 14 & 0 & exact \\
$e(\mathrm{lcm})\geq\max(e(a),e(\bar{a}))$ componentwise
  & 14 & 0 & exact \\
Bilateral support $\Leftrightarrow$ strict increase
  & 11 & 0 & exact \\
Degenerate pairs detected ($\Phi(a)\mid\Phi(\bar{a})$)
  & 3 & 0 & exact \\
\midrule
Divisibility preserves order (random pairs) & 10\,000 & 0 & exact \\
lcm $\geq$ max (random pairs) & 10\,000 & 0 & exact \\
\bottomrule
\end{tabularx}
\end{table}

Of the 11 non-degenerate duality pairs, 9 show strict increase on
$r$ (crystallised truth), and all~11 show strict increase on at
least one axis. The two pairs without $r$-increase
(\emph{orden}/\emph{caos}, \emph{creaci\'on}/\emph{destrucci\'on})
have interiors only in Layer~3 ($J$-axis), correctly predicted by
the bilateral support theorem. The two Layer~6 pairs
(\emph{temporal\_obs}/\emph{eterno\_obs},
\emph{receptivo}/\emph{creador\_obs}) show strict increase on all
four axes, as their interiors span all contributing layer sets.

\paragraph{What distinguishes dualities from arbitrary pairs.}
Duality pairs have a mean Jaccard overlap of $0.614$ (they share
most of their primitives), compared to $0.076$ for random pairs
of the same sizes (Kolmogorov--Smirnov $p < 10^{-4}$). This high
overlap is precisely what makes duality synthesis produce
\emph{small, targeted} gains: the mean complexity increase is
$1.43$ primes (vs.\ $5.97$ for random pairs, which share almost
nothing). The synthesis of opposites generates less raw complexity
than combining unrelated concepts, but the gain is
\emph{semantically structured}: it fills the specific gap between
the two poles, not arbitrary corners of the concept space.

%======================================================================
\section{Discussion}
\label{sec:discussion}
%======================================================================

\subsection{Measuring vs.\ enabling}

Constructive complexity (Section~\ref{sec:ccm}) and
constructor-theoretic information (Section~\ref{sec:ct}) address
complementary questions: \emph{how much was built?} (backward-looking)
and \emph{what can be built?} (forward-looking). The prime algebra
bridges both: it has a constructive complexity measure ($\omega$)
and satisfies constructor-theoretic conditions (I1--I5), suggesting
that these are two faces of the same coin.

\begin{theorem}[Complexity--constructor bridge]
\label{thm:bridge}
Let $(D,B,\oplus)$ be a domain with finite basis
$B=\{b_1,\ldots,b_k\}$, a join-composition $\oplus$ satisfying
$S(a\oplus b)=S(a)\cup S(b)$ for a support function
$S\colon D\to 2^B$, and unique decomposition (every element has a
unique support).
\begin{enumerate}[nosep]
\item[\textup{(i)}] $\mu(x)=|S(x)|-1$ is a CCM satisfying C1--C5.
\item[\textup{(ii)}] The medium satisfies I1--I4 of Constructor
  Theory.
\item[\textup{(iii)}] I5 (superinformation) is independent: it
  holds iff the medium admits a second-order gap
  $\mathrm{gap}(a,b)$ with $S(\mathrm{gap})=S(a)\setminus S(b)$.
\end{enumerate}
The prime algebra satisfies all three: FTA provides unique
decomposition, $\mathrm{lcm}$ is the join, and integer division
provides the gap.
\end{theorem}

\begin{proof}
\textup{(i)}
C1: $|S(b_i)|=1$, so $\mu(b_i)=0$.
C2: $|S(x)|\geq1$ for all $x$, so $\mu\geq0$.
C3: $|S(a)\cup S(b)|\leq|S(a)|+|S(b)|$, giving
$\mu(a\oplus b)\leq\mu(a)+\mu(b)+1$ with $c=1$.
C4: unique support over a finite basis is decidable (enumerate
and check).
C5: if $b\notin S(a)$, then $S(a\oplus b)=S(a)\cup\{b\}$,
so $\mu(a\oplus b)=\mu(a)+1>\mu(a)$.

\textup{(ii)}
I1: distinct elements have distinct supports by uniqueness;
the support serves as a distinguishing certificate.
I2: copying the support to a blank substrate clones the element.
I3: $|B|\geq2$ gives $\geq2$ distinguishable states.
I4: $\oplus$, support extraction, and equality testing are
computable over a finite basis.

\textup{(iii)}
C1--C5 constrain a function $\mu\colon D\to\mathbb{Z}_{\geq0}$
on individual elements. The gap takes two inputs and outputs a
relationship; by Proposition~\ref{prop:super}, it is irreducibly
second-order (varying the second argument changes the output even
when the first is fixed). Therefore I5 is not entailed by
C1--C5. Conversely, C1--C5 do not prohibit I5: the prime algebra
satisfies both.
\end{proof}

\subsection{Synthesis and the four axes}

The layer-based embedding of~\citet{p11} distributes the 72
primitives across four axis-contributing sets:
$R$ (Layers 1,2,4,6), $I$ (Layers 4,6), $J$ (Layers 3,5,6),
$K$ (Layer~6 only). This yields a precise connection between
synthesis and the quaternionic axes.

The Bilateral Synthesis Theorem
(Remark~\ref{rem:synth_coords}) states that axis $X$ shows
strict increase iff both interiors have $X$-support. Therefore:
\begin{itemize}[nosep]
\item A duality whose interiors differ only in Layer~3 or~5
  primitives (e.g., \emph{orden}/\allowbreak\emph{caos},
  \emph{creaci\'on}/\allowbreak\emph{destrucci\'on}) produces
  strict increase on $\exac{j}$ but not on $\rey{r}$: the synthesis
  generates new \emph{verifiable} structure.
\item A duality whose interiors include Layer~4 primitives
  (e.g., \emph{bien}/\emph{mal}, \emph{verdad}/\emph{mentira})
  produces strict increase on $\rey{r}$ and $\imag{i}$: the
  synthesis generates new crystallised truth \emph{and}
  potentiality.
\item Layer~6 dualities
  (\emph{temporal\_obs}/\emph{eterno\_obs},
  \emph{receptivo}/\emph{creador\_obs}) produce strict increase
  on all four axes, including $\recu{k}$: recursive selection
  is engaged only at the observer level.
\end{itemize}

This replaces the earlier informal mapping (blanket$\to r$,
interiors$\to i$, etc.)\ with a precise, layer-determined
characterisation: which axes increase depends on \emph{where} in
the dependency structure the interiors diverge, not on a fixed
assignment of synthesis components to axes.

\subsection{Relation to Lawvere's categorical synthesis}

The synthesis $a\oplus\bar{a}=\mathrm{lcm}(\Phi(a),\Phi(\bar{a}))$
shares the pattern of Lawvere's formalisation of Hegelian
\emph{Aufhebung}~\citep{lawvere1969,lawvere1996}: thesis and
antithesis produce a synthesis that preserves both. The precise
categorical status is as follows.

The prime algebra, ordered by divisibility, is a distributive
lattice and therefore a thin category. In this category,
$\mathrm{lcm}$ is the \emph{coproduct} (categorical join) and
$\gcd$ is the \emph{product} (categorical meet)---standard lattice
theory. The coprime-interiors theorem (Theorem~\ref{thm:mb}) adds
content beyond the coproduct: it guarantees that the join
decomposes uniquely into a shared blanket and disjoint
contributions from each pole.

Lawvere's Aufhebung, however, operates on \emph{levels} of a
cohesive topos via adjoint functors between connected-component
functors at different scales~\citep{lawvere1996}. Our algebra has
no topos structure, no connected-component functors, and no
adjunction between levels. The parallel is therefore at the level
of the \emph{dialectical pattern}---opposition $\to$ preservation
$\to$ elevation---not at the level of the categorical
formalism. We share Lawvere's conclusion (synthesis produces
genuinely new structure) but reach it by a different route
(unique factorisation in a distributive lattice, not adjoint
functors in a topos).

%======================================================================
\section{Conclusion}
\label{sec:conclusion}
%======================================================================

The imaginary dimensions of the quaternionic G-lattice are not
decoration. Through the synthesis of opposites---formalised as
$\mathrm{lcm}$ in the prime algebra---they generate compound
concepts whose constructive complexity strictly exceeds that of
either pole (Theorem~\ref{thm:complexity_increase}). The shared
structure between any two concepts partitions uniquely into an
algebraic Markov blanket and coprime interiors
(Theorem~\ref{thm:mb}), resolving the observer-dependence critique.
The gap operation---the raw material of synthesis---is irreducibly
second-order (Proposition~\ref{prop:super}), and three algebraic
tasks that synthesis enables are provably impossible under cosine
similarity (Theorem~\ref{thm:impossible}).

The Bilateral Synthesis Theorem of~\citet{p11} lifts these
results from prime factor counts to quaternionic coordinates,
characterising exactly which axes show strict increase for each
duality pair (Section~\ref{sec:computation}). Of 14 tested
dualities, the 11 non-degenerate pairs all confirm the theorem;
the 3 degenerate pairs (one pole subsuming the other) correctly
produce no increase. The prime algebra is shown to be irreducible
for this purpose: no function on quaternionic coordinates alone
can compute synthesis
(Proposition~\ref{prop:irreducible})---the embedding captures
geometric structure while the prime algebra retains combinatorial
identity. The bridge between constructive complexity (C1--C5) and
constructor-theoretic information (I1--I5) is made precise
(Theorem~\ref{thm:bridge}): for media with unique decomposition
the two frameworks are equivalent up to superinformation, which
is the irreducible second-order condition.

The companion paper~\citep{p11} showed that Boolean logic is what
survives when all perspectives are averaged: the context-free
special case. This paper shows what happens when perspectives
\emph{interact}: opposites, operating through the imaginary
dimensions, produce new truth. Boolean logic captures the panorama;
duality synthesis captures the engine.

%======================================================================
% References
%======================================================================
\bibliographystyle{plainnat}
\begin{thebibliography}{99}

\bibitem[Bruineberg et~al.(2018)]{bruineberg2018}
Bruineberg, J., Kiverstein, J. \& Rietveld, E. (2018).
The anticipating brain is not a scientist: the free-energy principle
from an ecological-enactive perspective.
\textit{Synthese}, 195, 2417--2444.

\bibitem[Deutsch(2013)]{deutsch2013}
Deutsch, D. (2013).
Constructor theory.
\textit{Synthese}, 190, 4331--4359.

\bibitem[Deutsch \& Marletto(2015)]{deutsch2015}
Deutsch, D. \& Marletto, C. (2015).
Constructor theory of information.
\textit{Proceedings of the Royal Society A}, 471, 20140540.

\bibitem[Lawvere(1969)]{lawvere1969}
Lawvere, F.W. (1969).
Adjointness in foundations.
\textit{Dialectica}, 23(3--4), 281--296.

\bibitem[Lawvere(1996)]{lawvere1996}
Lawvere, F.W. (1996).
Unity and identity of opposites in calculus and physics.
\textit{Applied Categorical Structures}, 4, 167--174.

\bibitem[Ornelas Brand(2026a)]{p1}
Ornelas Brand, J.A. (2026).
Triadic Neurosymbolic Engine: Prime Factorization as a Neurosymbolic
Bridge: Projecting Continuous Embeddings into Discrete Algebraic Space
for Deterministic Verification.
\textit{Zenodo}.
\href{https://doi.org/10.5281/zenodo.19205804}{doi:10.5281/zenodo.19205804}

\bibitem[Ornelas Brand(2026b)]{p4}
Ornelas Brand, J.A. (2026).
Toward a Layered Ontology of Semantic Duality: 14 Candidate Dualities
Across 6 Algebraic Layers with Preliminary Domain and Neural Evidence
(v1.0).
\textit{Zenodo}.
\href{https://doi.org/10.5281/zenodo.19375166}{doi:10.5281/zenodo.19375166}

\bibitem[Ornelas Brand(2026c)]{p11}
Ornelas Brand, J.A. (2026).
Quaternionic Logic: A G-Lattice Unifying Boolean and Fuzzy Frameworks (v0.1.0).
\textit{Zenodo}.
\href{https://doi.org/10.5281/zenodo.19560986}{doi:10.5281/zenodo.19560986}

\bibitem[Sharma et~al.(2023)]{sharma2023}
Sharma, A. et~al. (2023).
Assembly theory explains and quantifies selection and evolution.
\textit{Nature}, 622, 321--328.

\end{thebibliography}

\end{document}
